\section{Related Works}\label{sec:ch-7-related_work}

Automatic evaluation of summarization has historically centered on $n$-gram overlap with reference summaries, most notably ROUGE~\cite{lin2004rouge}. While simple and reproducible, ROUGE rewards surface overlap rather than semantic equivalence and correlates only weakly with human judgments of quality~\cite{liu2008correlation,bhandari2020reevaluating}. Embedding-based metrics such as BERTScore~\cite{zhang2020bertscore} and MoverScore~\cite{zhao2019moverscore} address this by comparing contextualized representations, and learned metrics like BLEURT~\cite{sellam2020bleurt} train directly on human ratings. A parallel line of work targets specific aspects of summary quality, particularly factual consistency, with metrics including FactCC~\cite{kryscinski2020evaluating}, QAGS~\cite{wang2020asking}, QuestEval~\cite{scialom2021questeval}, and SummaC~\cite{laban2022summac}. More recently, large language models have been used directly as evaluators of summary quality and have been shown to outperform earlier automatic metrics on several aspects of summary evaluation~\cite{liu2023geval,fu2024gptscore,chiang2023llmeval}. These metrics share a common assumption that summary quality can be assessed independently of the reader, which is precisely the assumption we challenge.

Query-focused summarization (QFS) generates summaries tailored to a user-provided query rather than producing a generic synopsis~\cite{daume2006bayesian,dang2005duc}, with approaches ranging from early extractive methods that scored sentences by query relevance~\cite{wan2007manifold,otterbacher2009biased} to neural systems that generate query-conditioned abstractive summaries~\cite{baumel2018query,xu2020coarse,vig2022exploring}, supported by datasets such as QMSum~\cite{zhong2021qmsum} and AQuaMuSe~\cite{kulkarni2020aquamuse}. A parallel line of work conditions summarization on reader attributes, including controllable systems that specify length, style, or focus~\cite{fan2018controllable,he2022ctrlsum}, audience-adaptive approaches targeting populations such as children, experts, or laypeople~\cite{chandrasekaran2020overview,august2022paperplain,goldsack2022making,luo2022readability}, and persona-based prompting of large language models~\cite{deshpande2023toxicity,zhang2024personallm}. However, QFS evaluation has largely relied on reference-based metrics applied against query-specific summaries~\cite{zhong2021qmsum,vig2022exploring}; our work extends this framing by pairing each query with an explicit reader persona and interrogating whether existing metrics can detect query-conditioned differences in informational content.

A number of datasets have been released with human annotations of summary quality, enabling meta-evaluation of automatic metrics. SummEval~\cite{fabbri2021summeval} provides expert ratings of system outputs on CNN/DailyMail across coherence, consistency, fluency, and relevance. FRANK~\cite{pagnoni2021frank} contributes fine-grained factual error annotations, and~\textcite{tang-etal-2023-understanding} extends this with error annotations across additional summarizers and datasets. More recent benchmarks elicit preference judgments via pairwise comparison of system outputs~\cite{goyal2022news,liu2023benchmarking}. Across these resources, annotations are collected without reference to a particular reader's goals; our evaluation anchors judgments to a specified persona and query, enabling direct measurement of information satisfaction.
